\documentclass{article}
\usepackage{amsmath,amssymb}
\usepackage{xurl}
\usepackage[hidelinks]{hyperref}
% Shared commands for notes in docs/paper-gaps/.

\DeclareUnicodeCharacter{2080}{\ifmmode{}_{0}\else\textsubscript{0}\fi}
\DeclareUnicodeCharacter{2081}{\ifmmode{}_{1}\else\textsubscript{1}\fi}
\DeclareUnicodeCharacter{2082}{\ifmmode{}_{2}\else\textsubscript{2}\fi}
\DeclareUnicodeCharacter{2099}{\ifmmode{}_{n}\else\textsubscript{n}\fi}

\newcommand{\Lean}{\textsc{Lean}}
\ExplSyntaxOn
\NewDocumentCommand{\leanid}{m}{
  \ifmmode
    \textsf{\detokenize{#1}}
  \else
    \group_begin:
    \urlstyle{sf}
    \tl_set:Nn \l_tmpa_tl {#1}
    \tl_replace_all:Nnn \l_tmpa_tl {\_} {_}
    \exp_args:NV \path \l_tmpa_tl
    \group_end:
  \fi
}
\ExplSyntaxOff
\providecommand{\tr}{\operatorname{tr}}
\newcommand{\tnleanIssueBase}{https://github.com/LionSR/TNLean/issues/}
\newcommand{\tnleanPrBase}{https://github.com/LionSR/TNLean/pull/}
\newcommand{\tnleanDocBase}{https://sirui-lu.com/TNLean/docs/}
\newcommand{\ghissue}[1]{\href{\tnleanIssueBase#1}{GitHub issue \##1}}
\newcommand{\ghpr}[1]{\href{\tnleanPrBase#1}{PR \##1}}
\newcommand{\doclink}[1]{\href{\tnleanDocBase#1.html}{\path{#1}}}

% Dirac notation and the identity operator, as in blueprint/src/macros/common.tex.
% \providecommand keeps note-local definitions authoritative.
\usepackage{dsfont}
\providecommand{\ket}[1]{|#1\rangle}
\providecommand{\bra}[1]{\langle #1|}
\providecommand{\braket}[2]{\langle #1 | #2 \rangle}
\providecommand{\ketbra}[2]{|#1\rangle\!\langle #2|}
\providecommand{\Id}{\mathds{1}}
\providecommand{\one}{\mathds{1}}
\providecommand{\C}{\mathbb{C}}
\providecommand{\MN}[1]{M_{#1}(\C)}
\providecommand{\MD}{M_D(\C)}

% Verdict marker: \gapnote{<kind>}{<status>}. Kind is the policy
% classification (clarification, local-correction, scope-restriction,
% unfaithful, false-source, open-gap); status is open, wip (elimination
% underway), resolved, or historical. Machine-read by the site index;
% prints a verdict line.
\newcommand{\gapnote}[2]{\par\noindent\textbf{Verdict:} #1 (#2).\par}

\title{Preparation in Logarithmic Depth: the Periodic State Must Not Vanish}
\author{}
\date{2026-09-28}
\begin{document}
\maketitle
\gapnote{local-correction}{resolved}

\section{What the source states}
Let $A$ be a normal tensor of physical dimension $d$ and bond dimension $D$.
Equation~(1) of~\cite{Malz2024Preparation} states that the normalized
translation-invariant state $\ket{\phi_N}$ generated by $A$ on a ring of $N$
sites is prepared with error $\varepsilon$ by a local circuit of depth
$O(\log(N/\varepsilon))$. The Supplemental Material writes the state as
$\ket{\phi_N}=c_N^{-1}\sum_{i_1,\dots,i_N}\operatorname{Tr}(A^{i_1}\cdots A^{i_N})
\ket{i_1\cdots i_N}$, ``where $c_N>0$ is the normalization constant'', and its
proof of Theorem~1 begins ``Let $\{\ket{\phi_N}\}$ be a sequence of TI
normalized MPS on $N$ sites generated by a normal tensor $A$''. Thus
$c_N=\lVert\phi_N(A)\rVert$ with
$\phi_N(A)(s)=\operatorname{Tr}(A^{s_1}\cdots A^{s_N})$, and the source takes
$c_N>0$.

\section{The degenerate case}
The condition $c_N>0$, that is $\ket{\phi_N(A)}\neq0$, can fail for a normal
tensor at small $N$. Take $d=4$, $D=3$, $\zeta=e^{2\pi i/3}$, and
\begin{align}
    A^1 & =E_{12}, & A^2 & =E_{23}, & A^3 & =E_{31},
        & A^4 & =\operatorname{diag}(1,\zeta,\zeta^2).
    \notag
\end{align}
Every product of two of these matrices has zero trace: a product containing an
off-diagonal matrix unit is off-diagonal or zero, and
$\operatorname{Tr}\bigl((A^4)^2\bigr)=1+\zeta^2+\zeta^4=0$. Hence $\ket{\phi_2(A)}=0$. The words of length
three span $M_3(\C)$: the words with three off-diagonal letters give the
diagonal matrix units, those with one off-diagonal letter give nonzero
multiples of $E_{12},E_{23},E_{31}$, and those with two give nonzero multiples
of $E_{13},E_{21},E_{32}$. So $A$ is normal.

For such $N$ there is no normalized state. With the convention that the
normalization of the zero vector is zero, every unit vector has error $1$, and
the claim fails for $\varepsilon<1$.

\section{The correction adopted}
The formal statement assumes $\ket{\phi_N(A)}\neq0$, the source's $c_N>0$.\footnote{Formal
    declaration:
    \leanid{MPSPreparation.exists_isPreparedInDepth_le_log_of_mpvState_ne_zero}
    in \doclink{TNLean/MPS/Preparation/LogDepthPreparation}.}
For every normal tensor there is $N_0$ with $\ket{\phi_N(A)}\neq0$ for all
$N\ge N_0$: in the gauge of eq.~(5) of the source,
$c_N^2=\operatorname{Tr}\mathcal E_A^N$ converges to $1$; the source's proof of
Lemma~1$'$(i) uses the same identity, ``$c_N = \sqrt{\tr (E_A^N)}$''.\footnote{Formal declarations:
    \leanid{MPSPreparation.exists_mpvState_ne_zero_of_le} and the resulting
    statement without the assumption,
    \leanid{MPSPreparation.exists_isPreparedInDepth_le_log}.}
Equation~(1) is an asymptotic statement, and for $N\ge N_0$ the assumption is
not needed.

\bibliographystyle{alpha}
\bibliography{references}
\end{document}
